% Part: normal-modal-logic
% Chapter: syntax-and-semantics
% Section: introduction

\documentclass[../../../include/open-logic-section]{subfiles}

\begin{document}

\olfileid[pt]{nml}{syn}{int}

\olsection{Introdução}

A lógica modal trata de \emph{proposições modais} e das relações de
acarretamento entre elas. Alguns exemplos de proposições modais são:
\begin{enumerate}
\item É necessário que $2+2=4$.
\item É necessariamente possível que chova amanhã.
\item Se é necessariamente possível que~$!A$, então é possível
  que~$!A$.
\end{enumerate}
Possibilidade e necessidade não são as únicas modalidades: outros
conectivos unários também são classificados como modalidades, por
exemplo, ``deveria ser o caso que~$!A$'', ``será o caso que~$!A$'',
``Dana sabe que~$!A$'' ou ``Dana acredita que~$!A$''.

A lógica modal aparece pela primeira vez no \emph{De Interpretatione}
de Aristóteles: ele foi o primeiro a notar que a necessidade implica a
possibilidade, mas não o inverso; que possibilidade e necessidade são
interdefiníveis; que, se $!A \land !B$ é possivelmente verdadeira,
então $!A$ é possivelmente verdadeira e $!B$ é possivelmente verdadeira,
mas não o inverso; e que, se $!A \to !B$ é necessário, então, se $!A$
é necessário, $!B$ também o é.

A primeira abordagem moderna à lógica modal foi a obra de C.~I. Lewis,
que culminou em \emph{Symbolic Logic} (1932), de Lewis e Langford.
Lewis \& Langford estavam insatisfeitos com a representação da
implicação por meio do condicional material: $!A \lif !B$ é um
substituto insatisfatório para ``$!A$ implica $!B$''. Em vez disso,
propuseram caracterizar a implicação como ``Necessariamente, se $!A$,
então $!B$'', simbolizada por $!A \strictif !B$. Ao tentar distinguir
as diferentes propriedades, Lewis identificou cinco sistemas modais,
\Log{S1}, \ldots, \Log{S4}, \Log{S5}; os dois últimos ainda são usados.

A abordagem de Lewis e Langford era puramente \emph{sintática}: eles
identificaram axiomas e regras razoáveis e investigaram o que se podia
provar por esses meios. Uma abordagem semântica permaneceu fora de
alcance por muito tempo, até que Rudolf Carnap fez uma primeira tentativa
em \emph{Meaning and Necessity} (1947), usando a noção de \emph{descrição
  de estado}, O.s.\@, uma coleção de sentenças atômicas (aquelas que são
``verdadeiras'' nessa descrição de estado). Depois de estender a definição
de verdade a sentenças arbitrárias $!A$, Carnap define $!A$ como
\emph{necessariamente verdadeira} se ela é verdadeira em todas as
descrições de estado. A abordagem de Carnap não podia tratar de
modalidades \emph{iteradas}, pois sentenças da forma ``Possivelmente,
necessariamente, \ldots, possivelmente $!A$'' sempre se reduzem à
modalidade mais interna.

O grande avanço da semântica modal veio com o artigo de Saul Kripke
``A Completeness Theorem in Modal Logic'' (JSL 1959). Kripke baseou
seu trabalho na ideia de Leibniz de que uma afirmação é necessariamente
verdadeira se é verdadeira ``em todos os mundos possíveis''. Essa ideia,
porém, sofre dos mesmos inconvenientes que a de Carnap, pois a verdade de
uma afirmação em um mundo $w$ (ou em uma descrição de estado $s$) não
depende de $w$. Kripke supôs, então, que os mundos estão relacionados por
uma \emph{relação de acessibilidade} $R$ e que uma afirmação da forma
``Necessariamente $!A$'' é verdadeira em um mundo $w$ se, e somente se,
$!A$ é verdadeira em todos os mundos $w'$ \emph{acessíveis a partir de}
$w$. Semânticas que oferecem alguma versão dessa abordagem são chamadas
semânticas de Kripke e tornaram possível o desenvolvimento tumultuado das
lógicas modais (no plural).

Quando interpretada pela semântica de Kripke, a lógica modal nos mostra
como são as \emph{estruturas relacionais} ``por dentro''. Uma estrutura
relacional é apenas um conjunto munido de uma relação binária (por
exemplo, o conjunto dos estudantes de uma turma ordenado pelo número de
seu cadastro de seguridade social é uma estrutura relacional). Na verdade,
porém, as estruturas relacionais aparecem em toda sorte de domínios: além
da possibilidade relativa de estados do mundo, podemos ter estados
epistêmicos de um agente relacionados pela possibilidade epistêmica, ou
estados de um sistema dinâmico com suas transições de estado, e assim por
diante. A lógica modal pode ser usada para modelar todos eles: o primeiro
caso nos dá a lógica modal alética ordinária; os demais nos dão a lógica
epistêmica, a lógica dinâmica etc.

Concentramo-nos em um ângulo específico, conhecido pelos lógicos modais
como ``teoria da correspondência''. Uma das primeiras descobertas mais
significativas de Kripke foi que muitas propriedades da relação de
acessibilidade~$R$ (se ela é transitiva, simétrica etc.) podem ser
caracterizadas \emph{na própria linguagem modal} por meio de ``esquemas
modais'' apropriados. Os lógicos modais dizem, por exemplo, que a
reflexividade de $R$ ``corresponde'' ao esquema ``Se necessariamente
~$!A$, então~$!A$''. Exploramos principalmente a teoria da correspondência
de vários sistemas clássicos de lógica modal (p.ex., \Log{S4} e
\Log{S5}) obtidos por uma combinação dos esquemas \Ax{D}, \Ax{T}, \Ax{B},
\Ax{4} e~\Ax{5}.

\end{document}
